% Part III -- Reusable and Reliable Programs
% Chapter 7: Useful Built-in Functions

\h{Useful Built-in Functions}

Python includes a set of functions that are ready to use without an import.
You have already met several of them, including \c{print()}, \c{len()},
\c{range()}, and \c{type()}. This chapter organizes the most useful built-ins
and shows when each one is a good choice.

\begin{NoteBox}
\textbf{Prerequisites.} You should understand collections, loops,
conditionals, and how to call a function.
\end{NoteBox}

\begin{tcolorbox}[title=Learning outcomes,colframe=orange,colback=white,arc=2mm]
After completing this chapter, you should be able to:
\begin{itemize}
    \item choose built-in functions for common numeric and collection tasks;
    \item distinguish functions that return a new value from methods that modify data;
    \item use \c{any()} and \c{all()} to summarize boolean values;
    \item use \c{enumerate()} and \c{zip()} clearly;
    \item inspect documentation with \c{help()}; and
    \item import focused tools from Python's standard library.
\end{itemize}
\end{tcolorbox}

\hh{Built-in means always available}

A built-in function belongs to Python itself. You call it directly without an
\c{import} statement.

\code{c07_familiar_builtins}{python}{
    room_names = ["Lobby", "Office", "Cafe"]

    print(len(room_names))
    print(type(room_names))
}{Call two familiar built-in functions}

The Python documentation lists every built-in, but beginners do not need to
memorize the full list. Learn a small set and look up the rest when needed.

\hh{Summarize numeric collections}

Use \c{sum()}, \c{min()}, and \c{max()} for common summaries.

\code{c07_numeric_summary}{python}{
    areas_m2 = [24.0, 35.5, 18.0, 31.0]

    total_m2 = sum(areas_m2)
    smallest_m2 = min(areas_m2)
    largest_m2 = max(areas_m2)
    average_m2 = total_m2 / len(areas_m2)

    print(f"Total: {total_m2:.1f} m2")
    print(f"Range: {smallest_m2:.1f} to {largest_m2:.1f} m2")
    print(f"Average: {average_m2:.1f} m2")
}{Summarize a list of areas}

These functions make the intent clearer than a manually written loop when only
a standard summary is needed.

\begin{NoteBox}
\textbf{Stop and predict.} What problem would \c{min([])} cause? An empty list
has no smallest or largest item. Chapter 8 explains how to handle such errors.
\end{NoteBox}

\hh{Work with magnitude and precision}

The \c{abs()} function removes the sign, and \c{round()} returns a value rounded
to a requested number of decimal places.

\code{c07_abs_round}{python}{
    difference_m = -0.1267

    magnitude_m = abs(difference_m)
    rounded_m = round(magnitude_m, 2)

    print(magnitude_m)
    print(rounded_m)
}{Use absolute value and rounding}

Use rounding for display or when the specification requires it. Keep full
precision during intermediate calculations when possible.

\hh{Sort without changing the original}

The \c{sorted()} function returns a new list. The original collection remains
unchanged.

\code{c07_sorted}{python}{
    floor_numbers = [3, 1, 4, 2]

    ascending = sorted(floor_numbers)
    descending = sorted(floor_numbers, reverse=True)

    print(floor_numbers)
    print(ascending)
    print(descending)
}{Create sorted lists without changing the original}

By contrast, the list method \c{floor_numbers.sort()} changes the list and
returns \c{None}. Choose based on whether the original order should be kept.

\code{c07_sort_method}{python}{
    floor_numbers = [3, 1, 4, 2]
    result = floor_numbers.sort()

    print(floor_numbers)
    print(result)
}{A sorting method modifies the list}

\hh{Reverse an order}

The \c{reversed()} function produces values in reverse order. Convert the
result to a list when you need to store or display all values at once.

\code{c07_reversed}{python}{
    room_names = ["Lobby", "Office", "Cafe"]
    reverse_order = list(reversed(room_names))

    print(reverse_order)
    print(room_names)
}{Reverse without modifying the original list}

\hh{Summarize boolean values}

Use \c{any()} when at least one \c{True} value is enough. Use \c{all()} when
every value must be \c{True}.

\code{c07_any_all}{python}{
    checks = [True, True, False, True]

    print(any(checks))
    print(all(checks))
}{Compare any with all}

The first result is \c{True} because at least one check passes. The second is
\c{False} because one check fails.

The boolean list can be created from data before it is summarized.

\code{c07_checks_from_values}{python}{
    widths_m = [1.2, 0.9, 1.5, 1.1]
    wide_enough = []

    for width_m in widths_m:
        wide_enough.append(width_m >= 1.0)

    print(wide_enough)
    print(all(wide_enough))
}{Build and summarize boolean checks}

\hh{Number and pair items}

You used \c{enumerate()} and \c{zip()} in Chapter 5. They are worth keeping in
your regular toolkit.

\code{c07_enumerate_zip}{python}{
    room_names = ["Lobby", "Office", "Cafe"]
    areas_m2 = [24.0, 18.0, 31.0]

    for number, pair in enumerate(zip(room_names, areas_m2), start=1):
        room_name, area_m2 = pair
        print(f"{number}. {room_name}: {area_m2:.1f} m2")
}{Number paired values}

When nested calls make a loop hard to read, create an intermediate variable or
use a simpler loop. Compact code is not automatically clearer code.

\hh{Convert between common types}

The names \c{int}, \c{float}, \c{str}, \c{list}, \c{tuple}, and \c{set} can be
called to create or convert values.

\code{c07_type_conversion}{python}{
    count_text = "4"
    count = int(count_text)
    area_text = str(24.5)
    letters = list("ABC")
    unique_levels = set([1, 2, 2, 3])

    print(count)
    print(area_text)
    print(letters)
    print(unique_levels)
}{Convert values and collections}

Conversion works only when the source value has a suitable form. For example,
\c{int("four")} raises an exception.

\hh{Check a type when it is genuinely needed}

The \c{isinstance()} function checks whether a value has a particular type.

\code{c07_isinstance}{python}{
    value = 3.5

    print(isinstance(value, float))
    print(isinstance(value, str))
}{Check a value's type}

Do not add type checks everywhere. Clear inputs and conversions are usually
easier to follow. Use \c{isinstance()} when a function intentionally accepts
more than one kind of value.

\hh{Read built-in documentation}

The \c{help()} function displays documentation in the REPL.

\code{c07_help_example}{python}{
    values = [3, 1, 2]

    print(sorted(values))
    # Try help(sorted) in the REPL.
}{Use help interactively when a function is unfamiliar}

Documentation commonly shows parameters inside square brackets when they are
optional. Try a small example after reading the description.

\hh{Import focused standard-library tools}

Not every useful function is built in. Python's standard library contains
modules that are installed with Python.

\code{c07_standard_library}{python}{
    import math
    from statistics import mean

    areas_m2 = [24.0, 35.5, 18.0, 31.0]
    radius_m = 3.0

    print(f"Mean area: {mean(areas_m2):.2f} m2")
    print(f"Circle area: {math.pi * radius_m ** 2:.2f} m2")
}{Use the math and statistics modules}

\c{math.pi} identifies where the name \c{pi} comes from. The targeted import
\c{from statistics import mean} allows a direct call to \c{mean()}.

\hh{Common mistakes}

\begin{itemize}
    \item \textbf{Expecting sorted() to change a list}: store its returned list.
    \item \textbf{Expecting list.sort() to return the sorted list}: it returns \c{None}.
    \item \textbf{Calling min() or max() on an empty collection}: check first.
    \item \textbf{Rounding too early}: keep precision until the result is displayed.
    \item \textbf{Using zip() with unequal lengths}: extra values are ignored.
    \item \textbf{Importing many names without purpose}: import only the tools you use.
\end{itemize}

\hh{Practice ladder}

\hhh{Level 0 -- Read and predict}

\code{c07_practice_level0}{python}{
    values = [6, 2, 9, 4]

    print(min(values))
    print(max(values))
    print(sum(values))
}{Predict three numeric summaries}

\hhh{Level 1 -- Modify}

Add two values to the list. Then display the average using \c{sum()} and
\c{len()}.

\hhh{Level 2 -- Complete}

Create ascending and descending versions of a list while preserving the
original.

\code{c07_practice_level2}{python}{
    levels = [4, 1, 3, 2]
    ascending = sorted(levels)
    descending = sorted(levels, reverse=True)

    print(levels)
    print(ascending)
    print(descending)
}{Preserve the original order}

\hhh{Level 3 -- Apply}

Given a list of areas, display the count, total, average, smallest value, and
largest value. Format decimal results to one decimal place.

\textbf{Hints}
\begin{itemize}
    \item \textbf{Hint 1}: Use \c{len()}, \c{sum()}, \c{min()}, and \c{max()}.
    \item \textbf{Hint 2}: Average is total divided by count.
    \item \textbf{Hint 3}: Store each result under a descriptive name.
\end{itemize}

\hhh{Level 4 -- Challenge}

Given room names and matching areas, display the rooms from largest area to
smallest. Hint: pair the data with \c{zip()}, then use \c{sorted()}.

\hh{Practice solutions}

\textbf{Level 0:} The outputs are \c{2}, \c{9}, and \c{21}.

\textbf{Level 2:} The original remains \c{[4, 1, 3, 2]}.

\textbf{Level 3}

\code{c07_solution_level3}{python}{
    areas_m2 = [24.0, 35.5, 18.0, 31.0]

    count = len(areas_m2)
    total_m2 = sum(areas_m2)
    average_m2 = total_m2 / count

    print(f"Count: {count}")
    print(f"Total: {total_m2:.1f} m2")
    print(f"Average: {average_m2:.1f} m2")
    print(f"Smallest: {min(areas_m2):.1f} m2")
    print(f"Largest: {max(areas_m2):.1f} m2")
}{One solution to Level 3}

\textbf{Level 4}

\code{c07_solution_level4}{python}{
    room_names = ["Lobby", "Office", "Cafe"]
    areas_m2 = [24.0, 18.0, 31.0]

    room_areas = list(zip(room_names, areas_m2))
    largest_first = sorted(room_areas, key=lambda pair: pair[1], reverse=True)

    for room_name, area_m2 in largest_first:
        print(f"{room_name}: {area_m2:.1f} m2")
}{One solution to Level 4}

The short \c{lambda} supplies a sort key. It is optional vocabulary; a named
function is better when the rule needs more explanation.

\hh{Chapter checkpoint}

\begin{enumerate}
    \item What makes a function built in?
    \item Which three functions provide a total, minimum, and maximum?
    \item Does \c{sorted()} modify its input list?
    \item What does \c{list.sort()} return?
    \item When does \c{any()} return \c{True}?
    \item When does \c{all()} return \c{True}?
    \item What risk comes with unequal collections passed to \c{zip()}?
    \item What is \c{help()} useful for?
    \item Why might you import \c{math} or \c{statistics}?
\end{enumerate}

\hh{Checkpoint answers}

\begin{enumerate}
    \item It is available without an import.
    \item \c{sum()}, \c{min()}, and \c{max()}.
    \item No. It returns a new list.
    \item \c{None}.
    \item When at least one supplied value is true.
    \item When every supplied value is true.
    \item Pairing stops when the shortest collection ends.
    \item Reading documentation for a function or object.
    \item They provide focused tools that are not built in.
\end{enumerate}

\begin{tcolorbox}[title=Chapter summary,colframe=orange,colback=white,arc=2mm]
You can select built-in functions for summaries, sorting, conversion, boolean
checks, pairing, and inspection. You can also import focused tools from the
standard library. The next chapter explains how to respond when an operation
cannot be completed.
\end{tcolorbox}

